\subsection{Transfer to Four Physical Event Families}
\label{sec:four-family-comparison}

Does a recovery forecast still help when a fault or a generator torque
pulse can also explain the PMUs? We tested a four-family version on
a reduced IEEE 39-bus plant with the same eight PMU sites. Its 31
coordinates comprise eight voltage magnitudes and angles, eight
current magnitudes and seven current-direction terms. It fits the
first fault or line opening from the measured voltages, propagates the
estimated generator and governor states after clearance, and compares
new reports with continued recovery and with a second intervention.
This implementation scores a bank of network responses with a Student
likelihood. The two-family filter in Section~\ref{sec:narrative-method}
uses a local Gaussian state update for each history. Both ask whether the
next transient requires a new switch. The four-family predictor uses
the reduced simulator's network equations. The selected
operating point was fixed before we generated 15 new scenario blocks.
The blocks reuse the reserved asset catalog and draw new times, event
strengths, noise and governor parameters. The supplement gives the fit,
the event prior and the alarm rule for this extension.

\begin{table}[t]\centering\footnotesize
\caption{Four-family extension on the reduced plant. Joint requires the correct family and asset. Extras count unmatched alarms. Quiet$^+$ counts quiet windows labeled abnormal.}
\label{tab:four-family-comparison}\setlength{\tabcolsep}{2.6pt}
\begin{tabular*}{\columnwidth}{@{\extracolsep{\fill}}lrrrr@{}}\toprule
\multicolumn{5}{l}{\textit{Continuous monitoring, 420 events in 465 records}}\\
Method & Macro-F1 & Family & Joint & Extras \\
\midrule
Recovery bank & 0.976 & 405 & 399 & 3 \\
With voltage check & 0.958 & 396 & 391 & 9 \\
Earlier physical bank & 0.957 & 396 & 387 & 10 \\
Sequential GLR & 0.835 & 326 & 321 & 16 \\
Refitted ExtraTrees & 0.708 & 265 & 252 & 0 \\
\midrule
\multicolumn{5}{l}{\textit{Family at a common alarm, 120 scheduled second events}}\\
Readout & Correct & Wrong & Late & Missed \\
\midrule
Recovery bank & 106 & 1 & 7 & 6 \\
CNN-F & 28 & 79 & 7 & 6 \\
CNN-F, no subtraction & 27 & 80 & 7 & 6 \\
MultiRocket-Hydra & 50 & 57 & 7 & 6 \\
\midrule
\multicolumn{5}{l}{\textit{One decision at 30 s, 72 isolated events and 12 quiet records}}\\
Method & Macro-F1 & Family & Joint & Quiet$^+$ \\
\midrule
Recovery bank & 0.983 & 69 & 69 & 0 \\
Archived ExtraTrees & 0.067 & 2 & 0 & 0 \\
\bottomrule\end{tabular*}\end{table}


The continuous test in Table~\ref{tab:four-family-comparison} counts a
joint decision when the correct family and exact bus or line are
published within 1.5~s of onset. The recovery bank met this deadline in
\RepBankJoint{} of 420 events (95\% block-bootstrap interval
$\RepBankJointCI$\%) and named \RepBankActive{} while the intervention
was still active. The refitted trees made \RepTreeJoint{}
joint diagnoses. The two methods issued \RepBankExtras{} and
\RepTreeExtras{} unmatched alarms, respectively. Compared with the earlier
physical bank, the recovery bank added 12 joint diagnoses within the
deadline and both made \RepBankActive{} before clearance. Ten missed
generator pulses and four generator-bus errors accounted for 14 of its
21 unsuccessful joint diagnoses. The extra voltage check discarded
eight correct joint diagnoses on balance and increased unmatched alarms
from \RepBankExtras{} to \RepCheckExtras{}.

For a first generator pulse, we asked which of the existing terminals
could resolve a bus that voltages alone leave ambiguous. We propagated
a seven-report torque pulse through each generator and let every rival
bus adjust its onset, duration, inertia and amplitude. At the 1.5-s
deadline, buses 35 and 36 remained the closest pair. The current angle
at bus 22 on branch 21--22 gave the largest separation among seven
available terminals. We froze two first-event readers. One used the
eight voltage phasors, and the other added this current angle. Both
summed over possible pulse onsets before choosing a bus. On a further
15 blocks the original bank, voltage reader and terminal reader
identified \TorqueBaseJoint{}, \TorqueVoltJoint{} and
\TorqueTermJoint{} of 420 events within 1.5~s, respectively. All ten
additional joint diagnoses over the bank were isolated short pulses.
The terminal reader located \TorqueTermShort{} of these 30 pulses,
including \TorqueTermShortActive{} while torque was still applied.
The bank located \TorqueBaseShort{}, including \TorqueBaseShortActive{}
while active. At bus 35, adding terminal 22 raised the voltage reader's
timely diagnoses from \TorqueVoltBusThirtyFive{} to
\TorqueTermBusThirtyFive{} of ten. The original bank issued
\TorqueBaseExtra{} unmatched alarms and the terminal version issued
\TorqueTermExtra{}. Both additional alarms correctly named bus 32 one
report after the deadline in the two governor variants of the same
scenario. We charged both as unmatched. The supplement reports the
paired intervals and sensitivity to terminal-angle noise.

The two late declarations motivated a revision. For a first generator
pulse we project the voltage reports and all seven current angles
onto the rotor-motion subspace and integrate over torque magnitude.
During recovery we check whether the shunt implied by a new fault
declaration would produce the observed voltage jump. We froze both
changes and repeated the experiment on 15 new scenario blocks from
the same reserved assets. The revised estimator triggered on 416 of
420 events, correctly classified all 416, and located 409 by the
unchanged 1.5-s deadline (97.38\%, 95\% block-bootstrap interval
95.71--98.81\%). The single-current estimator located 408. The
additional diagnosis was a short pulse at bus 35. Across all 30
short pulses, the revision raised the number located before clearance
from 10 to 16. Both estimators had macro-F1 0.9896 and made the
same five unmatched declarations, all on streams with two line trips.

We also compared family decisions on the same alarm reports. Each
classifier published 29 reports after the alarm when that report was
available.
All 120 scheduled second events count, including those whose trigger
was absent or too late. The spectral CNN-F uses the subtraction of
Li and Wang \cite{liwang2019successive}. MultiRocket-Hydra uses
convolutional time-series features and a ridge classifier
\cite{dempster2023hydra,middlehurst2024bakeoff}. Both learned four families on the same
400 simulated training cases. The second events in this test were
generator, load or line interventions. The bank retained its family
decision from the alarm and reported it at the common deadline. The
recovery bank, CNN-F and MultiRocket-Hydra correctly labeled
\RepBankSecond{}, \RepCnnSecond{} and \RepHydraSecond{} second events,
respectively. Equal publication times leave a difference in what each
model remembers. The bank carries a network state from earlier reports,
CNN-F can extrapolate a previous disturbance, and MultiRocket-Hydra
receives the two seconds around the alarm. Six missed triggers and
seven late readouts cost every method the same 13 second events.

Finally, the archived ExtraTrees system assigns one family and
location to a complete window. We gave it 84 independently simulated
30-s windows at its native duration. At window close we compared its
label with the bank's first declaration or normal if it remained
silent. Quiet$^+$ counts quiet records labeled as events. This measures
transfer of the archived weights to the reduced plant. Event strengths,
the balanced three-phase reconstruction and the measurement noise all
belong to this simulation.
